\begin{proposition}[Extension of the marked torus period]\hypertarget{tag-01GD}{}\label{cdp:st:D.2.6}
\StatementRecord{cdp:st:D.2.6}{proved}{01GD}{\SeriesRef{cdp:st:D.2.3}, \SeriesRef{cdp:eq:D.20}}


Let $h:U\to\Delta$ be a proper holomorphic map from a smooth
complex threefold, smooth with connected two-torus fibres over
$\Delta^*$. Suppose that its integral $H^1$ monodromy is the
identity. The marked period map
\[
 \Delta^*\longrightarrow
 \mathcal T_2
 :=\{W\in\operatorname{Gr}(2,H^1(U_t;\C)):
                   W\oplus\overline W=H^1(U_t;\C)\}
\]
extends holomorphically across zero with value in $\mathcal T_2$.
In particular the relative Albanese group family over $\Delta^*$
extends to a smooth proper group family of complex two-tori over
$\Delta$. 
\end{proposition}


\begin{proof}

\begin{enumerate}[leftmargin=*]
\item The punctured family is connected by local triviality and connectedness
of its fibres, so it lies in one open and closed component $U^0$ of $U$.
The inclusion $U^0\hookrightarrow U$ is an isomorphism over $\Delta^*$.
Apply analytic embedded resolution to the compact reduced central divisor
in $U^0$ \cite[\S0.2(2), preprint p.~3; Theorem~1.6 and Remarks~1.7,
preprint pp.~5--6]{BierstoneMilman1997}. This gives a proper composite
of blowups $\pi:\widetilde U\to U^0$ with smooth centres over that
divisor, with $\widetilde U$ smooth and the reduced central support
of $\widetilde h:=h|_{U^0}\circ\pi$ normal crossing. For $t\ne0$,
$\pi_t:\widetilde U_t\to U_t$ is a biholomorphism;
its pullback identifies holomorphic forms, the integral marking and
cup products. Transporting through these pullbacks, use the relative
logarithmic complex in the proof of Proposition~\ref{cdp:st:D.2.3}
on this model, with its full central multiplicities.
Identity on $H^1$ gives identity on $H^2=\bigwedge^2H^1$.
Consequently its logarithmic Gauss--Manin bundle has zero
residue and a holomorphic flat frame. Let $\Omega$ be a
generator of $\widetilde h_*\mathscr K^2$, as in Proposition~\ref{cdp:st:D.2.3}.
For $t\ne0$, let $\Omega_t$ be the form on $U_t$ determined by
$\pi_t^*\Omega_t=\Omega|_{\widetilde U_t}$.

\item The coefficient-ideal argument in
\eqref{cdp:eq:D.18}--\eqref{cdp:eq:D.19} shows that $\Omega$
cannot vanish on the whole reduced central support: otherwise
its divided cohomology class would have bounded periods while
its positive $L^2$ volume diverges. Thus $\Omega|_T$ is nonzero
on some normalized component $T$. The annular volume estimate
in the same proof makes every conductor residue of every
component restriction zero: a nonzero residue would give a
divergent positive volume while its cup-product expression
stays bounded. The nonzero restriction $\Omega|_T$ is therefore
a holomorphic two-form on $T$.

\item Choose a point of $T$ outside the conductor and the zeros of
this two-form. In a small chart the map is $t=z^m$ and
\[
 \Omega=a(z,\xi,\eta)\,d\xi\wedge d\eta,
 \qquad a(0,0,0)\ne0.
\]
Integrating on one sheet over a fixed small
$(\xi,\eta)$-polydisc yields a constant $c>0$ such that
\begin{equation}
 \int_{U_t}\Omega_t\wedge\overline{\Omega_t}\ge c
 \qquad(0<|t|\ll1).
\hypertarget{tag-0197}{}\label{cdp:eq:D.20}
\end{equation}
The cohomology class $v(t)=[\Omega_t]$, expressed in the flat
frame, is holomorphic across zero. Its cup pairing with its conjugate equals the
integral in \eqref{cdp:eq:D.20}, hence
\[
 v(0)\wedge\overline{v(0)}>0
 \quad\hbox{in }H^4(U_t;\R)\simeq\R,
\]
using the complex orientation. In particular $v(0)\ne0$.

\item Write $V=H^1(U_t;\C)$ with its fixed integral marking.
For $t\ne0$, $v(t)$ is a decomposable element of
$\bigwedge^2V$, spanning $\bigwedge^2H^{1,0}(U_t)$.
The decomposable cone is the closed algebraic set defined by
the Pl\"ucker equations, so $v(0)$ is decomposable. For a
nonzero decomposable $v$, the exact inverse to the Pl\"ucker
map is
\[
 W(v)=\{\alpha\in V:\alpha\wedge v=0\}.
\]
This kernel has dimension two and varies holomorphically near
$v(0)$. The nonzero pairing
$v(0)\wedge\overline{v(0)}$ says exactly that
$W(v(0))\cap\overline{W(v(0))}=0$. The extended period
therefore lies in $\mathcal T_2$.

\item Let $\mathcal W\subset V\otimes\OO_\Delta$ be the extended
rank-two subbundle $W(v(t))$, and let
$\Lambda=H_1(U_t;\Z)$ with its fixed marking. Evaluation of
cohomology on integral cycles gives a holomorphic embedding
of the constant lattice $\Lambda$ in $\mathcal W^*$.
A real cycle annihilating $\mathcal W_t$ also annihilates
$\overline{\mathcal W_t}$, hence all of $V$; it is therefore zero.
Equality of real dimensions makes the real-linear evaluation
map an isomorphism, so the image is a full lattice in every fibre,
uniformly on compact subsets of the disc. Translation by this
lattice is free and properly discontinuous. For a fixed integral
basis, the closed parallelotopes generated by its images cover
each quotient fibre; over a compact subset of the base their
union is compact. Consequently
\[
 \mathcal A=\mathcal W^*/\Lambda\longrightarrow\Delta
\]
is a smooth proper family of complex Lie groups with
two-torus fibres. On $\Delta^*$, integration of holomorphic
one-forms identifies each group fibre with the Albanese torus
of $U_t$, and these identifications vary holomorphically.


\end{enumerate}
\end{proof}